\subsection{Proof of the main theorem}
\label{subsec:technical-main-proof}

\begin{table}[tbp]
\centering\small
\renewcommand{\arraystretch}{1.3}
\begin{tabular}{@{}p{0.20\textwidth}p{0.45\textwidth}p{0.28\textwidth}@{}}
\toprule
\textbf{Parameter} & \textbf{Explicit choice} & \textbf{Scaling in $r$}\\
\midrule
$M$ & $2^{27}$ & Fixed\\
$\gamma$ & $1-10^{-6}$ & Fixed\\
$C_1$ & $2\times10^8$ & Fixed\\
$C_{2,r}$ & $10001\sqrt{2}\sqrt{(M+9)^r-M^r}$ & $\Theta((M+9)^{r/2})$\\
$L_r$ & $\left\lceil\dfrac{\log(2P_r/E)}{2\log\theta}\right\rceil$ & $\Theta(rM^{2r})$\\[3pt]
$Q_r$ & $2^{q_r}$, with $q_r$ in \eqref{eq:q-choice} & $\exp(\Theta(r^2M^{2r}))$\\
Input dimension $n_r$ & $k^2Q_r^{2r}$ & $\exp(\Theta(r^3M^{2r}))$\\
Output dimension $k$ & $M^r$ & $\exp(\Theta(r))$\\
\bottomrule
\end{tabular}
\caption{Parameters of the final channel $\widehat\Lambda_r$.
The basic channel $\Phi_r$ has input dimension $Q_r^{2r}$.
The quantities $E,\theta,P_r$ are defined below.}
\label{tab:all-r}
\end{table}

We use the parameters in Table~\ref{tab:all-r}. Thus
$M=2^{27}$, $\gamma=1-10^{-6}$, and $C_1=2\times10^8$.
Throughout this proof, $r\ge1$ is an integer. Put
\[
 k=M^r,\qquad A_r=(M+9)^r-M^r,\qquad
 E=10^{-6},\qquad \theta=\frac{10001}{10000},\qquad
 P_r=\left(2\left\lceil\sqrt{C_1}k\right\rceil\right)^{k(k-1)}.
\]
Choose
\begin{equation}\label{eq:q-choice}
 \begin{gathered}
 C_{2,r}=10001\sqrt{2A_r},\qquad
 L_r=\left\lceil\frac{\log(2P_r/E)}{2\log\theta}\right\rceil,\\
 Q_r=2^{q_r},\qquad
 q_r=2+\left\lceil L_r\bigl[4\log_2(4M-1)+2\log_2 k\bigr]\right\rceil.
 \end{gathered}
\end{equation}
These choices satisfy $C_{2,r}>\sqrt{C_1A_r}$, $L_r\ge1$,
and give $P_r\theta^{-2L_r}\le E/2$ and
$Q_r>B_{L_r}k^{2L_r}$. Since $A_r\ge1$, Proposition~\ref{prop:app-two-copy} yields
\begin{equation}\label{eq:main-filter-bound}
 \frac{\operatorname{Tr}F}{Q_r^{2r}}
 \le P_r\theta^{-2L_r}
 \left[1+\frac{B_{L_r}}{Q_r}
 \left(\frac{k(k-1)}{A_r}\right)^{L_r}\right]<E.
\end{equation}
Let $\Phi_r$ be the resulting basic channel and write
\begin{equation}\label{eq:main-a}
 a_r=1+\frac{\gamma^2}{k}
 \left(\frac{C_{2,r}}{\sqrt{C_1}-\sqrt2}\right)^2.
\end{equation}
The matrices $U_I$ are real, and the grid $\mathcal W_k$ is closed
under entrywise conjugation. Hence $F$ and $H=\sqrt{\gamma}(I+F)^{-1/2}$ are real, so $\overline{\Phi_r}=\Phi_r$.
Proposition~\ref{prop:measurement-support}, Lemma~\ref{lemma:app-single-copy}, and Corollary~\ref{cor:supp-tensor-entropy} give
\begin{equation}\label{eq:main-basic-bounds}
 k\left\|\Phi_r(\sigma)-\frac{I_k}{k}\right\|_2^2\le a_r-1,
 \qquad
 H\bigl((\Phi_r\otimes\Phi_r)(\omega_{Q_r^{2r}})\bigr)
 \le 2\log k-rd_M\!\left(\frac{\gamma^2}{(1+E)^2}\right),
\end{equation}
where $\omega_{Q_r^{2r}}$ is the maximally entangled input state defined above.

We next make every one-copy output close to the maximally mixed state
using a short average of Pauli conjugations. Set
\[
 m=27r,\qquad
 t_r=\left\lceil\log_2(54r^2\sqrt{a_r})\right\rceil,
 \qquad h_r=2^{t_r},\qquad \ell_r=h_r^2.
\]
Represent $\mathbb F_{h_r}$ using the lexicographically first monic
irreducible binary polynomial of degree $t_r$, and let
$\tau:\mathbb F_{h_r}\to\mathbb F_2$ be its field trace.
For $x,y\in\mathbb F_{h_r}$, define
\[
 P_{x,y}=\bigotimes_{j=1}^{m}
 X^{\tau(yx^{j-1})}Z^{\tau(yx^{m+j-1})},
 \qquad
 \mathcal R_r(Y)=\frac1{\ell_r}\sum_{x,y}P_{x,y}YP_{x,y}^{\dagger}.
\]
Here $X,Z$ are the qubit Pauli matrices and $x^0=1$.
For every nonidentity element $P$ of the Pauli basis, conjugation contributes a
sign $(-1)^{\tau(yp(x))}$ for a nonzero binary polynomial $p$ of degree
at most $2m-1$. Averaging first over $y$ gives
\[
 \mathcal R_r(P)=\frac{\#\{x:p(x)=0\}}{h_r}P,
 \qquad
 \frac{\#\{x:p(x)=0\}}{h_r}\le\frac{2m-1}{h_r}
 <\frac1{r\sqrt{a_r}}.
\]
Orthogonality of the Pauli basis and the definition of $t_r$ imply
\begin{equation}\label{eq:main-randomizer}
 \|\mathcal R_r(Y)\|_2\le\frac{\|Y\|_2}{r\sqrt{a_r}}
 \quad(\operatorname{Tr}Y=0),\qquad
 \ell_r<108^2r^4a_r.
\end{equation}
Consequently $\Lambda_r=\mathcal R_r\circ\Phi_r$ satisfies
\[
 k\left\|\Lambda_r(\sigma)-\frac{I_k}{k}\right\|_2^2
 \le\frac{a_r-1}{r^2a_r}\le r^{-2}.
\]
The inequality $H(\rho)\ge-\log\operatorname{Tr}(\rho^2)$ now gives
\begin{equation}\label{eq:main-smoothed-single}
 H_{\min}(\Lambda_r)\ge\log k-\log(1+r^{-2}).
\end{equation}
On the other hand, $(\mathcal R_r\otimes\mathcal R_r)$ averages
$\ell_r^2$ unitary conjugations. Applying the entropy-of-mixture
inequality to the Bell output in \eqref{eq:main-basic-bounds} gives
\begin{equation}\label{eq:main-smoothed-double}
 H\bigl((\Lambda_r\otimes\Lambda_r)(\omega_{Q_r^{2r}})\bigr)
 \le 2\log k-rd_M\!\left(\frac{\gamma^2}{(1+E)^2}\right)
       +2\log\ell_r.
\end{equation}

We obtain the final channel by appending a classical label.
For $u,v\in\{0,1\}^{m}$, write
$W_{uv}=\bigotimes_{j=1}^{m}X^{u_j}Z^{v_j}$, and define
\begin{equation}\label{eq:main-weyl-extension}
 \widehat\Lambda_r(\rho)
 =\sum_{u,v}W_{uv}
 \Lambda_r\!\left(\langle u,v|\rho|u,v\rangle\right)W_{uv}^{\dagger}.
\end{equation}
Its input and output dimensions are $k^2Q_r^{2r}$ and $k$.
Every output is a mixture of unitary conjugates of outputs of $\Lambda_r$.
Entropy concavity therefore gives
\[
 \chi_{\max}(\widehat\Lambda_r)\le\log k-H_{\min}(\Lambda_r).
\]
For two copies, use the Bell input and choose both labels independently
and uniformly. The average output is $I_{k^2}/k^2$, while all states
in this ensemble have the same entropy. Hence
\[
 \chi_{\max}(\widehat\Lambda_r^{\otimes2})
 \ge 2\log k-H\bigl((\Lambda_r\otimes\Lambda_r)(\omega_{Q_r^{2r}})\bigr)
 \ge rd_M\!\left(\frac{\gamma^2}{(1+E)^2}\right)-2\log\ell_r.
\]

It remains to check that the two-copy gain exceeds the cost of this
short average. Put $\beta=(10001\gamma/9999)^2$. Then
$1<\beta<1.001$, and the binomial theorem gives
\[
 a_r=1+\beta\left[\left(1+\frac9M\right)^r-1\right]
 \le\left(1+\frac{9\beta}{M}\right)^r.
\]
For $s=\gamma^2/(1+E)^2$, the elementary estimate
\[
 d_M(s)\ge\frac{s\log M+s\log s+(1-s)\log(1-s)}{M}
\]
and $-s\log s\le1-s$ imply
\begin{align*}
 d_M(s)-2\log\left(1+\frac{9\beta}{M}\right)
 &\ge\frac{\log M-(1-s)[\log M+1-\log(1-s)]-18\beta}{M}\\
 &>\frac{0.692868}{M}>5\times10^{-9}.
\end{align*}
Here we used only $18.711<\log M<19$,
$1-s<4\times10^{-6}$, $\log(1-s)>-13$, and $\beta<1.001$.
The first estimate follows by writing $1+(M-1)s=M(s+(1-s)/M)$
and using $(x+y)\log(x+y)\ge x\log x+y\log y$ for $x,y>0$.
Combining the preceding bounds proves
\begin{equation}\label{eq:main-final-bounds}
 \chi_{\max}(\widehat\Lambda_r)\le\log(1+r^{-2}),\qquad
 \chi_{\max}(\widehat\Lambda_r^{\otimes2})
 >5\times10^{-9}r-8\log r-4\log108.
\end{equation}

The sharper asymptotic coefficient follows from
$a_r=1+\beta[(1+9/M)^r-1]\le\beta(1+9/M)^r$:
\begin{equation}\label{eq:main-sharp-holevo}
 \chi_{\max}(\widehat\Lambda_r^{\otimes2})
 \ge r\left[d_M(s)-2\log\left(1+\frac9M\right)\right]
       -8\log r-4\log108-2\log\beta.
\end{equation}
The corresponding upper bound holds for the entropy of the actual Bell
output of $\Lambda_r^{\otimes2}$, by \eqref{eq:main-smoothed-double}.
Thus the remainder in the stated asymptotic estimate is explicitly
controlled by $8\log r+4\log108+2\log\beta$.

Given $N\ge1$, take $J=\lceil N\rceil$ and $r_N=10^{11}J$, and define
$\Phi_N=\widehat\Lambda_{r_N}$. Since $\log(1+x)<x$ for $x>0$, the first bound in
\eqref{eq:main-final-bounds} is less than $r_N^{-2}<1/N$.
Using $\log(10^{11})<26$, $4\log108<20$, and $\log J\le J-1$,
the second bound is greater than
\[
 500J-8(26+\log J)-20\ge492J-220>N.
\]
Also $C(\Phi_N)\ge\tfrac12\chi_{\max}(\Phi_N^{\otimes2})$, so
$C(\Phi_N)-\chi_{\max}(\Phi_N)>N/2-1/N$.
Finally, $\log P_r=\Theta(rM^{2r})$ gives the scalings in
Table~\ref{tab:all-r}. Since $r_N=\Theta(N)$, the output dimension
is $\exp(\Theta(N))$ and the input dimension is
$\exp(\exp(\Theta(N)))$, as claimed.
